\section*{第 \ref{ch:g4:changes-of-state} 章——熔化与沸腾：状态变化}

\begin{solution}{exo:g4:changes-of-state:1}
\omterm{def:g2:ice-water-steam:meltfreeze}{熔化}；\omterm{def:g2:ice-water-steam:meltfreeze}{凝固}；沸腾；液化。
\end{solution}

\begin{solution}{exo:g4:changes-of-state:2}
\omterm{def:g4:changes-of-state:melting-point}{熔点}为 \qty{0}{\celsius}，\omterm{def:g4:changes-of-state:boiling-point}{沸点}为 \qty{100}{\celsius}。
巧克力在 \qty{34}{\celsius} 附近\omterm{def:g2:ice-water-steam:meltfreeze}{熔化}——接近口腔\omterm{def:g3:temperature-thermometers:temperature}{温度}。
\end{solution}

\begin{solution}{exo:g4:changes-of-state:3}
\qty{0}{\celsius}——五分钟后仍是 \qty{0}{\celsius}，
因为冰与水共存时，\omterm{def:g3:temperature-thermometers:temperature}{温度}被
固定在\omterm{def:g4:changes-of-state:melting-point}{熔点}上。
\end{solution}

\begin{solution}{exo:g4:changes-of-state:4}
物质在\omterm{def:g2:ice-water-steam:meltfreeze}{熔化}或沸腾时，\omterm{def:g3:temperature-thermometers:temperature}{温度}在
这一“里程碑”处保持不变。热量全部用于状态变化本身——
即\omterm{def:g3:solids-liquids-gases:liquid}{液体}变为水蒸气——而不用于升温。
\end{solution}

\begin{solution}{exo:g4:changes-of-state:5}
在 \qty{20}{\celsius} 时为\omterm{def:g3:solids-liquids-gases:solid}{固体}（低于 $60$）；在
\qty{80}{\celsius} 时为\omterm{def:g3:solids-liquids-gases:liquid}{液体}（高于 $60$）；在 \qty{35}{\celsius} 时仍为\omterm{def:g3:solids-liquids-gases:solid}{固体}，
——尽管口袋里的巧克力条在此\omterm{def:g3:temperature-thermometers:temperature}{温度}下早已\omterm{def:g2:ice-water-steam:meltfreeze}{熔化}。
\end{solution}

\begin{solution}{exo:g4:changes-of-state:6}
\omterm{def:g3:temperature-thermometers:temperature}{温度}并不更高：两锅都保持在 \qty{100}{\celsius}，
即沸腾时的平台。多余的热量只是让水蒸气产生得更快——火力白白浪费，
煮出的意面并无两样。
\end{solution}

\begin{solution}{exo:g4:changes-of-state:7}
依据曲线上的两段平坦部分，即平台。凡是曲线在
炉子持续加热时保持水平，便说明
正在发生状态变化。
\end{solution}

\begin{solution}{exo:g4:changes-of-state:8}
金（$1064 < 1200$）可以倾倒出来；铁（$1538 > 1200$）仍为\omterm{def:g3:solids-liquids-gases:solid}{固体}。炉温
未能达到铁的\omterm{def:g4:changes-of-state:melting-point}{熔点}，相差 $1538 - 1200 = 338$ 度。
\end{solution}

\begin{solution}{exo:g4:changes-of-state:9}
冰在\omterm{def:g2:ice-water-steam:meltfreeze}{熔化}时把饮料的\omterm{def:g3:temperature-thermometers:temperature}{温度}固定在 \qty{0}{\celsius}：流入的热量
全部用于\omterm{def:g2:ice-water-steam:meltfreeze}{熔化}冰，没有一点用于给饮料升温。只有当最后一点冰
也化完，\omterm{def:g3:temperature-thermometers:temperature}{温度}才开始上升。
\end{solution}

\begin{solution}{exo:g4:changes-of-state:10}
她的最佳猜测是：这种物质是冰，即冻结的水。所依据的规律是：
\omterm{def:g4:changes-of-state:melting-point}{熔点}和\omterm{def:g4:changes-of-state:boiling-point}{沸点}是物质的特征标记；纯净物质的\omterm{def:g4:changes-of-state:melting-point}{熔点}与\omterm{def:g4:changes-of-state:boiling-point}{沸点}
永不改变，故与水的两个数值完全吻合，
即可作为可靠的判据。
\end{solution}

\begin{solution}{exo:g4:changes-of-state:11}
人们一直确信水总在 \qty{100}{\celsius} 时沸腾。在高山上，
水的\omterm{def:g4:changes-of-state:boiling-point}{沸点}远低于此——水很热，却不足
\qty{100}{\celsius}，因此这“沸腾”的水\omterm{def:g3:temperature-thermometers:temperature}{温度}太低，不能
正确冲泡茶叶。罪魁祸首是较稀薄的\omterm{def:g2:air-around-us:air}{空气}对
水面的压力较小，这将在气压一章中说明。
\end{solution}
